\chapter{§4.2 平均数、中位数与众数}
\label{sec:4.2}


\prerequisite{§1.5 分数, §1.6 小数}

\gradelevel{5-7 年级}


一组数据往往包含很多个数值。我们能不能用\textbf{一个数}来代表这组数据的"中心"位置？这就是\textbf{集中趋势}的问题。本章介绍三种最常用的集中趋势度量：平均数、中位数和众数。



\section*{平均数（5 年级）}


\begin{explore}


学校组织植树活动，五个小组种树的棵数分别是：

\end{explore}

\begin{quote}
第 1 组：8 棵，第 2 组：6 棵，第 3 组：10 棵，第 4 组：7 棵，第 5 组：9 棵。
\end{quote}


老师问："平均每组种了多少棵？"这个问题该怎么理解？

\begin{understand}


\textbf{直觉理解}：如果把所有树重新平均分给每个小组，每组能分到几棵？

\[
\text{总棵数} = 8 + 6 + 10 + 7 + 9 = 40 \text{（棵）}
\]


\[
\text{平均每组} = 40 \div 5 = 8 \text{（棵）}
\]


\textbf{定义}：一组数据的\textbf{平均数}（也叫\textbf{算术平均数}），等于所有数据的总和除以数据的个数。

设 $n$ 个数据为 $x_1, x_2, \ldots, x_n$ ，则平均数：

\[
\bar{x} = \frac{x_1 + x_2 + \cdots + x_n}{n} = \frac{1}{n}\sum_{i=1}^{n} x_i
\]


\textbf{注意}：平均数不一定等于数据中的某个实际值。例如上面的例子中，平均数是 8，而实际数据中只有第 1 组恰好是 8。

\end{understand}

\begin{keytakeaway}


\begin{itemize}
  \item 平均数 $= \frac{\text{数据总和}}{\text{数据个数}}$ 。
  \item 平均数利用了每一个数据的信息，因此\textbf{每个数据都会影响平均数}。
  \item 平均数是最常用的集中趋势度量。
\end{itemize}


\end{keytakeaway}



\section*{加权平均数（6-7 年级）}


\begin{explore}


小华期末总评成绩的计算方法是：平时成绩占 $30\%$ ，期中考试占 $30\%$ ，期末考试占 $40\%$ 。他的三项成绩分别是 85 分、78 分、92 分。

如果直接算普通平均数： $(85 + 78 + 92) \div 3 = 85$ 分。

但这没有考虑到各项成绩的\textbf{重要程度不同}——期末考试占比最大，应该对总评的影响也更大。

\end{explore}

\begin{understand}


当各数据的\textbf{重要程度不同}时，需要计算\textbf{加权平均数}。每个数据乘以它的\textbf{权重}（即重要程度所占的比例），然后求和。

\[
\bar{x}_{\text{加权}} = x_1 \times w_1 + x_2 \times w_2 + \cdots + x_n \times w_n
\]


其中 $w_1 + w_2 + \cdots + w_n = 1$ （权重之和为 1）。

小华的总评：

\[
\bar{x} = 85 \times 0.3 + 78 \times 0.3 + 92 \times 0.4 = 25.5 + 23.4 + 36.8 = 85.7 \text{（分）}
\]


可以看到，加权平均数（ $85.7$ 分）和普通平均数（ $85$ 分）不一样。因为期末考试的权重最大，而小华期末考得最好，所以加权平均数更高。

\textbf{另一种常见的加权情境}——频率作为权重：

如果一组数据中某些值重复出现，也可以用加权平均数来简化计算。例如，10 名同学的跳远成绩中，3 人跳了 $3.2$ 米，5 人跳了 $3.5$ 米，2 人跳了 $3.8$ 米，则：

\[
\bar{x} = 3.2 \times \frac{3}{10} + 3.5 \times \frac{5}{10} + 3.8 \times \frac{2}{10} = 0.96 + 1.75 + 0.76 = 3.47 \text{（米）}
\]


\end{understand}

\begin{workedexamples}


\textbf{例 1.} 一家公司招聘，笔试成绩占 $40\%$ ，面试成绩占 $60\%$ 。小李笔试 80 分，面试 90 分；小王笔试 95 分，面试 75 分。谁的总评更高？

\textbf{解：}

\[
\text{小李总评} = 80 \times 0.4 + 90 \times 0.6 = 32 + 54 = 86 \text{（分）}
\]


\[
\text{小王总评} = 95 \times 0.4 + 75 \times 0.6 = 38 + 45 = 83 \text{（分）}
\]


$86 > 83$ ，所以小李的总评更高。虽然小王笔试比小李高 15 分，但面试权重更大（ $60\%$ ），面试小李远高于小王，因此小李总评胜出。

\end{workedexamples}

\begin{keytakeaway}


\begin{itemize}
  \item 当各数据重要程度不同时，使用加权平均数。
  \item 权重越大的数据，对加权平均数的影响越大。
  \item 普通平均数是加权平均数的特例——所有权重相等。
\end{itemize}


\end{keytakeaway}



\section*{中位数（7 年级）}


\begin{explore}


某公司有 7 名员工，月薪（单位：元）分别是：

\end{explore}

\begin{quote}
3000, 3500, 4000, 4200, 4500, 5000, 50000
\end{quote}


计算平均数：

\[
\bar{x} = \frac{3000 + 3500 + 4000 + 4200 + 4500 + 5000 + 50000}{7} = \frac{74200}{7} \approx 10600 \text{（元）}
\]


平均月薪约 10600 元！但 7 个人中有 6 个人的月薪都低于 10600 元。这个"平均数"能代表这家公司大多数人的薪资水平吗？

\textbf{问题出在哪里？} 50000 是一个\textbf{极端值}（远大于其他数据），它把平均数"拉"高了。我们需要一种不受极端值影响的度量。

\begin{understand}


\textbf{中位数}是将数据\textbf{从小到大排列}后，处于\textbf{正中间位置}的那个数。

求法：
\begin{enumerate}
  \item 将 $n$ 个数据从小到大排列。
  \item 如果 $n$ 是\textbf{奇数}，中位数是第 $\frac{n+1}{2}$ 个数据。
  \item 如果 $n$ 是\textbf{偶数}，中位数是第 $\frac{n}{2}$ 个与第 $\frac{n}{2}+1$ 个数据的平均值。
\end{enumerate}


\textbf{回到上面的例子}：数据已按从小到大排列，共 7 个（奇数），中位数是第 $\frac{7+1}{2} = 4$ 个数据，即 $4200$ 元。

$4200$ 元比 $10600$ 元更能代表大多数员工的薪资水平。

\textbf{偶数个数据的例子}：如果有 8 个数据 $2, 5, 7, 8, 10, 12, 15, 20$ ，则中位数为第 4 个和第 5 个数据的平均值：

\[
\text{中位数} = \frac{8 + 10}{2} = 9
\]


\end{understand}

\begin{workedexamples}


\textbf{例 2.} 10 名同学的数学考试成绩为：72, 85, 90, 65, 78, 92, 88, 76, 95, 82。求这组数据的中位数。

\textbf{解：}

第一步，从小到大排列： $65, 72, 76, 78, 82, 85, 88, 90, 92, 95$ 。

第二步， $n = 10$ （偶数），中位数是第 5 个和第 6 个数据的平均值。

第 5 个数据 $= 82$ ，第 6 个数据 $= 85$ 。

\[
\text{中位数} = \frac{82 + 85}{2} = 83.5 \text{（分）}
\]


\textbf{例 3.} 在例 2 中，如果新增一名同学的成绩为 80 分，中位数变为多少？

\textbf{解：}

新的排列： $65, 72, 76, 78, 80, 82, 85, 88, 90, 92, 95$ ，共 $11$ 个（奇数）。

中位数是第 $\frac{11+1}{2} = 6$ 个数据 $= 82$ 分。

\end{workedexamples}

\begin{keytakeaway}


\begin{itemize}
  \item 中位数是排列后处于正中间位置的值。
  \item 中位数\textbf{不受极端值影响}，这是它相对于平均数的最大优势。
  \item 求中位数前，必须先将数据\textbf{从小到大排列}。
\end{itemize}


\end{keytakeaway}



\section*{众数（7 年级）}


\begin{explore}


一家鞋店想进货一批运动鞋。为了决定各号码的进货量，店员记录了一天中卖出的鞋子的号码：

\end{explore}

\begin{quote}
38, 39, 40, 40, 40, 41, 41, 42, 40, 39, 40, 41, 40, 38, 40
\end{quote}


你认为鞋店应该重点进哪个号码的货？

\begin{understand}


这个问题中，平均数和中位数都不是最有用的。鞋店真正关心的是\textbf{哪个号码卖得最多}。

\textbf{众数}就是一组数据中\textbf{出现次数最多}的那个值。

统计各号码出现的次数：

\begin{center}
\begin{tabular}{llllll}
\toprule
号码 & 38 & 39 & 40 & 41 & 42 \\
\midrule
次数 & 2 & 2 & 7 & 3 & 1 \\
\bottomrule
\end{tabular}
\end{center}


40 号出现了 7 次，最多，所以众数是 $40$ 。鞋店应该重点进 40 号的货。

\textbf{特别说明}：
\begin{itemize}
  \item 一组数据可以有\textbf{多个众数}（如果有两个或更多值出现次数相同且最多）。
  \item 也可能\textbf{没有众数}（如果所有值出现次数都一样）。
\end{itemize}


\end{understand}

\begin{keytakeaway}


\begin{itemize}
  \item 众数是出现次数最多的值。
  \item 众数最适合回答"哪个最常见"的问题。
  \item 众数不需要计算，只需要\textbf{数频次}。
\end{itemize}


\end{keytakeaway}



\section*{三种统计量的比较与选择（7 年级）}


\begin{explore}


某班 9 名同学参加公益活动的捐款金额（单位：元）为：

\end{explore}

\begin{quote}
10, 10, 10, 20, 20, 30, 50, 50, 200
\end{quote}


算一算：
\begin{itemize}
  \item 平均数 $= \frac{10 \times 3 + 20 \times 2 + 30 + 50 \times 2 + 200}{9} = \frac{30 + 40 + 30 + 100 + 200}{9} = \frac{400}{9} \approx 44.4$ 元
  \item 中位数 $= 20$ 元（第 5 个数据）
  \item 众数 $= 10$ 元（出现 3 次，最多）
\end{itemize}


三种统计量给出了三个不同的数值！到底该用哪个来描述这组数据？

\begin{understand}


\begin{center}
\begin{tabular}{llll}
\toprule
统计量 & 含义 & 优势 & 劣势 \\
\midrule
平均数 & 所有数据的"均匀化"值 & 利用了全部数据信息 & 受极端值影响大 \\
中位数 & 排在正中间的值 & 不受极端值影响 & 没有利用全部数据的具体值 \\
众数 & 出现最多的值 & 反映"最常见"的情况 & 可能不唯一或不存在；不能反映整体水平 \\
\bottomrule
\end{tabular}
\end{center}


\textbf{选择的原则}：

\begin{itemize}
  \item 当数据分布比较均匀、没有极端值时，\textbf{平均数}最常用。
  \item 当数据中有明显的极端值（如薪资数据、房价数据）时，\textbf{中位数}更合适。
  \item 当关心"出现最多的是哪个"时（如衣服号码、投票结果），用\textbf{众数}。
  \item 在许多实际问题中，可以同时报告多个统计量，给出更完整的信息。
\end{itemize}


\end{understand}

\begin{workedexamples}


\textbf{例 4.} 两个射击运动员在 10 次练习中的成绩（环数）如下：

甲：7, 8, 7, 9, 8, 10, 7, 8, 8, 8

乙：5, 10, 10, 6, 8, 10, 9, 4, 10, 8

（1）分别求甲、乙的平均数、中位数和众数。
（2）你认为谁的成绩更好？请说明理由。

\textbf{解：}

（1）\textbf{甲}：

排列： $7, 7, 7, 8, 8, 8, 8, 8, 9, 10$

\[
\bar{x}_{\text{甲}} = \frac{7 \times 3 + 8 \times 5 + 9 + 10}{10} = \frac{21 + 40 + 9 + 10}{10} = \frac{80}{10} = 8 \text{（环）}
\]


中位数： $n = 10$ ，第 5 个 $= 8$ ，第 6 个 $= 8$ ，中位数 $= \frac{8+8}{2} = 8$ 环。

众数： $8$ （出现 5 次）。

\textbf{乙}：

排列： $4, 5, 6, 8, 8, 9, 10, 10, 10, 10$

\[
\bar{x}_{\text{乙}} = \frac{4 + 5 + 6 + 8 \times 2 + 9 + 10 \times 4}{10} = \frac{4 + 5 + 6 + 16 + 9 + 40}{10} = \frac{80}{10} = 8 \text{（环）}
\]


中位数：第 5 个 $= 8$ ，第 6 个 $= 9$ ，中位数 $= \frac{8+9}{2} = 8.5$ 环。

众数： $10$ （出现 4 次）。

（2）甲和乙的平均数都是 $8$ 环，似乎不相上下。但仔细观察：

\begin{itemize}
  \item 甲的成绩集中在 $7 \sim 10$ 之间，比较稳定。
  \item 乙的成绩从 $4$ 到 $10$ 波动很大，虽然高分更多，但也有很低的分数。
\end{itemize}


如果看\textbf{稳定性}，甲更好（数据波动小）。如果看\textbf{高分能力}，乙的众数 $10$ 说明他更常打出满分。这告诉我们：\textbf{只看一种统计量可能不够全面}，有时需要综合考虑。（在 §4.3 中我们将学习用"方差"来精确衡量数据的波动程度。）

\end{workedexamples}

\begin{keytakeaway}


\begin{itemize}
  \item 平均数、中位数、众数分别从不同角度描述数据的"中心"。
  \item 没有哪个统计量是万能的，选择取决于\textbf{数据特征}和\textbf{分析目的}。
  \item 当数据有极端值时，中位数通常比平均数更有代表性。
\end{itemize}


\end{keytakeaway}

\begin{practice}


\textbf{1.} 某超市一周内某款饮料每天的销量（瓶）为： $45, 52, 38, 52, 60, 52, 41$ 。
（1）求这组数据的平均数、中位数和众数。
（2）如果超市想知道每天应准备多少瓶该饮料，你建议用哪个统计量？

\textbf{2.} 某班 8 名同学 1 分钟仰卧起坐的个数为： $25, 30, 28, 32, 45, 29, 31, 28$ 。
（1）求平均数和中位数。
（2）如果有一名同学的成绩改为 $80$ 个（特别突出），平均数和中位数各变为多少？
（3）哪个统计量受极端值的影响更小？

\textbf{3.} 一次面试中，综合成绩由笔试（权重 $0.3$ ）、面试表现（权重 $0.5$ ）和简历评分（权重 $0.2$ ）三部分组成。张三的三项得分分别为 88 分、92 分、75 分，李四的三项得分分别为 95 分、80 分、90 分。谁的综合成绩更高？

\end{practice}



\begin{answer}


\textbf{1.}

（1）排列： $38, 41, 45, 52, 52, 52, 60$ 。

平均数 $= \frac{38 + 41 + 45 + 52 + 52 + 52 + 60}{7} = \frac{340}{7} \approx 48.6$ 瓶。

中位数： $n = 7$ （奇数），第 4 个数据 $= 52$ 瓶。

众数： $52$ （出现 3 次）。

（2）建议用\textbf{众数} $52$ 或\textbf{中位数} $52$ 。因为超市关心的是"大多数时候卖多少"，而平均数受到某天低销量（ $38$ ）的影响偏低，可能导致准备不足。

\textbf{2.}

（1）排列： $25, 28, 28, 29, 30, 31, 32, 45$ 。

平均数 $= \frac{25 + 28 + 28 + 29 + 30 + 31 + 32 + 45}{8} = \frac{248}{8} = 31$ 个。

中位数： $n = 8$ （偶数），第 4 个 $= 29$ ，第 5 个 $= 30$ ，中位数 $= \frac{29 + 30}{2} = 29.5$ 个。

（2）将 $45$ 改为 $80$ ，新排列： $25, 28, 28, 29, 30, 31, 32, 80$ 。

新平均数 $= \frac{25 + 28 + 28 + 29 + 30 + 31 + 32 + 80}{8} = \frac{283}{8} = 35.375$ 个。

中位数不变：第 4 个 $= 29$ ，第 5 个 $= 30$ ，中位数 $= 29.5$ 个。

（3）中位数不受极端值影响。平均数从 $31$ 变为 $35.375$ ，变化明显；中位数仍为 $29.5$ ，没有变化。

\textbf{3.}

\[
\text{张三} = 88 \times 0.3 + 92 \times 0.5 + 75 \times 0.2 = 26.4 + 46 + 15 = 87.4 \text{（分）}
\]


\[
\text{李四} = 95 \times 0.3 + 80 \times 0.5 + 90 \times 0.2 = 28.5 + 40 + 18 = 86.5 \text{（分）}
\]


张三的综合成绩（ $87.4$ 分）更高。虽然李四笔试和简历分都更高，但面试权重最大（ $0.5$ ），张三面试高出 12 分，最终胜出。

\end{answer}
